\section{Chapter 17}

\begin{verse}
The people at the beginning hardly knew that there were (sovereigns)\\
Their successors (i.e., sovereigns) were loved and exalted\\
Their successors were feared\\
Their successors were despised\\
Their perfidy lacking loyalty destroyed all trust\\
The first, solemn, reserved in speaking\\
Fulfilled their function perfectly\\
And the ten thousand beings said:\\
``We are living according to our nature.''
\end{verse}

\paragraph{Commentary}


Up until recent times in China the ``invisibility of the Emperor'' (because, by custom, he rarely showed himself alone) was an outer symbol of the invisible ``not-acting'' government. The text indicates the successive forms, always more degraded, assumed by the type of the leader.

\begin{enumerate}


\item First of all ``distance'' fails (the first principle was: ``he exercised an influence, only keeping himself distant''). The sovereign is visible and ``popular'': he loves popularity and needs it. This prestige is not based on the feeling of distance, but on that of closeness to him of the governed mass; for this reason he is loved.

\item In the next case the type of the prince follows, who governs only because he is feared.

\item We finally reach the leader who is feared and despised at the same time, failing every relationship of loyalism and mutual fidelity, and the structure of a State reaches the limit of instability.
\end{enumerate}


The sovereign of the origins, ``whose power derives from that of the Principle'', was called ``The Mysterious One'' --- says Chuang Tze. He transmits the influence of the Principle to all beings, making their natural capacities develop in them; at the origin his ``politics'' consisted ``in leading the individual nature of being back to conformity with the universal acting virtue.'' (XII, 1-10) The general idea is summarized rather well by Granet (La pensée chinoise, p 547) in the following terms:

\begin{quotex}
For the benefit of all things, but without charity nor pride, he is limited to concentrate in himself an intact Majesty ... this sovereign Majesty is not distinguished from pure Power nor from integral knowledge (\emph{i she}) ... An unknown Autocrat, he does his work without anyone noticing it, and this work is accomplished without affecting in any way the one from whom it emanates.
\end{quotex}


\paragraph{Chinese text and literal translation}


Chapter 17 (\begin{CJK*}{UTF8}{bsmi}第十七章\end{CJK*})

\columnratio{.43}
\begin{paracol}{2}
\begin{verse}
\begin{CJK*}{UTF8}{bsmi}
太上，下知有之；\\
其次，親而譽之；\\
其次畏之；\\
其次侮之。\\
信不足焉，有不信焉。\\
悠兮，其貴言，\\
功成事遂，\\
百姓皆謂我自然。
\end{CJK*}
\end{verse}
\switchcolumn
\begin{verse}
Great rulers are hardly known by their subjects,\\
Then come those the people draw near and praise,\\
Then those the people hold in fear,\\
Then those the people revile.\\
When one lacks trust, one finds no trust.\\
Reluctantly, without boasting;\\
Perform actions, accomplish deeds;\\
The people will say it happened naturally.
\end{verse}
\end{paracol}


\flrightit{Posted on 2021-02-05 by Lao Tzu}

\begin{center}* * *\end{center}

\begin{footnotesize}
\begin{sffamily}
\texttt{roballenvalecourt on 2021-02-05 at 23:09 said:}

I am constantly impressed with the Sicily connection in regards to Evola and the Western Esoteric Path... \end{sffamily}
\end{footnotesize}

